% Ad Familiares 16.15 --- headnote
% ad-familiares-16-15, 12--13 April 53 BC, Cumae

\textit{Cicero to Marcus Tullius Tiro, written from the villa at
Cumae on the day before the Ides of April --- 12 April 53 BC ---
with a postscript dashed off the same evening or the next morning,
13 April. Tiro is still ill, stranded somewhere on the road back
to his patron, and Cicero is now two letters deep in an anxious
vigil: 16.14 had gone out the day before, and this one tracks
the next two messengers as they straggle in. Aegypta arrives
first with the good news that the fever has broken; then Hermia
finally appears with a letter in Tiro's own hand, the writing
unsteady from weakness.}

\textit{The texture is unguarded and almost domestic. Cicero
counts up the messengers, names two slaves (Aegypta and a cook
he is sending back to keep house for Tiro on the road), and
falls into a small chiasmus that is the heart of the note:
\textit{qua si me liberaris, ego te omni cura liberabo} ---
``if you free me from that worry, I will free you from every
care.'' \textit{Vacillantibus litterulis} --- ``a tottering
little hand'' --- is the kind of physical detail he never wastes
on bigger correspondents; this is the prose of someone who has
been reading Tiro's drafts for twenty years and can see the
illness in the letter shapes.}
